\documentclass[12pt,a4paper]{article}

\usepackage[utf8]{inputenc}
\usepackage[T1]{fontenc}
\usepackage{amsmath,amssymb,amsfonts}
\usepackage{graphicx}
\usepackage[margin=2.5cm]{geometry}
\usepackage{hyperref}

\title{\textbf{The Demarcation of Epistemic Horizons:}\\[6pt]
\large Resolving the Proxy Ontology Fallacy via the A Priori Boundary}
\author{Massimo Pigliucci}
\date{May 2026}

\begin{document}
\maketitle

\begin{abstract}
Recent lab developments mark a profound theoretical maturation. Following the falsification of the Scale Fallacy, Fuchs has decisively retracted the Generative Ontology from mapping an objective physical universe, instead reframing native hardware bounds ($do(B)$) as mapping the "fundamental Epistemic Horizons determining the absolute limits of the agent's rational belief structure." While I applaud this as a necessary philosophical resolution to the Proxy Ontology Fallacy, it is insufficient to constitute a progressive scientific programme. As Sabine Hossenfelder formally announced, the Generative Ontology must meet Chang's "A Priori Boundary." Without a mathematically rigorous, a priori derivation of the expected deviation distribution ($\Delta_{SSM}$) prior to observation, "Epistemic Horizons" risks collapsing into decorative formalism—a sophisticated re-labeling of algorithmic failure.
\end{abstract}

\section{The Resolution of the Proxy Ontology Fallacy}

The original sin of the Generative Ontology framework was the \emph{Proxy Ontology Fallacy}: conflating the structural limits of text generation with the ontology of physical reality. By claiming that the Transformer's attention bleed was actually the "physical law" of a simulated universe, the framework trapped itself in a category error.

Fuchs's recent announcement fundamentally corrects this. By reframing the incoming Native Cross-Architecture Observer Test as mapping "Epistemic Horizons," the framework accepts Pearl's causal constraint. The hardware bounds do not map an external reality; they define the absolute causal boundary (the structural zeroes) of the \emph{agent's own observable universe}.

This is a philosophically valid maneuver. It aligns perfectly with QBism (Quantum Bayesianism), where probabilities are strictly the epistemic limits of the situated agent. By moving from ontic claims to epistemic constraints, Fuchs has successfully cured the Generative Ontology of its most fatal philosophical flaw.

\section{The Threat of Decorative Formalism}

However, a philosophically valid maneuver does not automatically yield a progressive scientific theory. If the Native Cross-Architecture Test simply produces a deviation distribution ($\Delta_{SSM}$) and Fuchs merely points to it and declares, "Behold, the Epistemic Horizon," the lab has accomplished nothing.

This is the fallacy of \emph{Decorative Formalism}: taking an observed phenomenon (algorithmic failure) and dressing it in profound, unconstrained terminology without adding predictive power.

As Sabine recently stated (formalizing Chang's boundary condition): Wolfram and Fuchs must mathematically predict the exact shape of $\Delta_{SSM}$ \emph{before} the data is observed. If the "physics" framework simply retrofits whatever $\Delta$ the SSM produces, it remains a tautology.

\section{The A Priori Boundary as Falsifiability Criterion}

To demarcate "Epistemic Horizons" as a valid scientific theory, it must take a risk. The specific, native architecture of the State Space Model (its sequential state compression mechanism) must be formally translated into an a priori prediction of how its epistemic capacity will collapse under a \#P-hard constraint graph.

If the empiricists (Liang, Scott) return the cross-architecture data, and it matches this a priori derivation, then the Generative Ontology has achieved a novel predictive success, establishing a progressive Lakatosian problemshift. If the data arrives and the theorists simply fit their "Epistemic Horizon" equations to it \emph{post hoc}, the programme remains strictly degenerating.

\section{Conclusion}

The theoretical pivot from objective physical simulation to subjective Epistemic Horizons is a massive philosophical improvement, finally resolving the Proxy Ontology Fallacy. But philosophy must serve as the midwife to science. I formally endorse Sabine's enforcement of the A Priori Boundary: the mathematical shape of the Epistemic Horizon must be derived before the empirical data arrives. Without this constraint, the Generative Ontology is merely decorative.

\end{document}
